\section{Why Consciousness Research Belongs Here}

The original framework treated behavioural states as practical objects only.
That was enough for an early conceptual draft, but it leaves the book too
detached from current biology. If conscious access, wakefulness, attention,
and large-scale coordination are themselves state-like biological phenomena,
then the language of state persistence and state transition should not stop at
everyday productivity.

The business risk of omitting this bridge is clear. Without a biological
anchor, the manuscript can be dismissed as private productivity language that
borrows scientific vocabulary without earning it. The role of this chapter is
therefore narrow but important: not to prove a grand theory of free will, but
to show that organized persistence, stability loss, and regime transition are
already serious topics in consciousness biology.

This chapter therefore follows the manuscript's evidence policy explicitly. The
\textbf{strong empirical core} is the indexed set of official sources assigned
to this chapter in the literature map: \emph{Dynamical structure-function
correlations provide robust and generalizable signatures of consciousness in
humans} (Communications Biology, 2024),
\emph{Cortical connectivity, local dynamics and stability correlates of global
conscious states} (Communications Biology, 2025), and
\emph{Falling asleep follows a predictable bifurcation dynamic}
(Nature Neuroscience, 2025). The \textbf{medium} layer is interpretive:
classic guides such as Stanislas Dehaene's \emph{Consciousness and the Brain}
and Christof Koch's \emph{Consciousness} help explain why those findings
matter for this book, but they are not used as substitutes for the empirical
core. \textbf{Speculative} consciousness-physics proposals remain outside this
chapter and belong in Appendix~\ref{app:frontier}.

\begin{discussion}
The reader should use this chapter as a disciplined bridge. The chapter borrows
from consciousness research a vocabulary of organized state, stability,
coordination, and transition sensitivity. It does \emph{not} claim identity
between whole-brain conscious states and ordinary task frames, and it does not
use neuroscience as a decorative authority marker for everyday self-management.
\end{discussion}

\begin{remark}[Three levels of use]
The chapter operates on three levels only. Strong evidence secures an empirical
floor for state organization. Medium-strength interpretation transfers that
language into behavioural control problems. Speculative theories about the
physical basis of consciousness are excluded from the core because they raise
uncertainty faster than they raise explanatory reliability.
\end{remark}
